% Synchronized editorial extracts from main.tex, 2026-10-07 pre-submission audit.
% main.tex is canonical; this file is not included in the build.

\begin{abstract}
Human-to-robot transfer combines perception, retargeting, calibration, and execution timing, making task failures difficult to attribute. We introduce the \emph{Transfer Gap Budget}, a paired oracle-intervention framework that estimates each component's recoverable contribution to success and selects targeted corrections. Across UR5e and UR10e simulations, perception dominates under the nominal error model, recovering approximately 30 percentage points. However, setting timing perturbations to span speeds measured in six human demonstrations makes timing dominant. The diagnosis is thus conditional on the transfer-error distribution. On human-video demonstrations, diagnosis-driven corrections and a revised retreat improve success on unused seeds; the notebook example rises from 0/60 to 60/60 on both arms. These findings motivate measuring demonstration-associated errors when choosing which transfer component to improve. All robot evaluations are in simulation.
\end{abstract}

\section{Introduction}

Human demonstrations provide motion and timing, but transferring them requires perception, motion retargeting, tool-frame calibration, and executable timing. Mismatches at several stages can produce the same task failure. Success alone therefore does not reveal which component should be improved; embodiment adaptation remains a broader challenge for robot learning~\citep{domae2026embodimentgap}.

As human demonstration datasets grow in quantity and diversity, it matters whether additional data addresses the dominant transfer bottleneck, a question relevant to Scaling H2R~\citep{scalingh2r}. Different demonstrations can change motion and timing statistics, while varied scenes can change perception challenges. The resulting error distribution may alter which component limits success. We study this diagnostic question using six demonstrations from three people to ground timing ranges and provide motion for simulated transfer.

We introduce the \emph{Transfer Gap Budget}: paired oracle substitutions measure each component's recoverable success contribution and select a correction on development seeds. We show that perception dominates under the nominal error model, but demonstration-derived timing makes timing dominant. Diagnosis-driven corrections and a retreat revision improve human-video transfer on unused seeds. The framework diagnoses which transfer component limits the use of human data.

\section{Limitations and Conclusion}
\label{sec:limits}

The budget is conditional on the assumed or measured error distribution. Timing ranges come from six demonstrations and depth-error magnitudes from simulation; calibration and retargeting errors and correction models remain assumed. The study uses one simulated task, scripted controllers, and a UR5e-tuned layout; layout shifts leave the table and task fixed. Three demonstrators and unreconstructed scenes limit generalisation. Look-again and probing manoeuvres and their time costs are not simulated. Real-robot performance is untested.

The Transfer Gap Budget diagnoses recoverable transfer failures through paired interventions. Perception dominates under the nominal error model; demonstration-derived timing makes timing dominant. Corrections and a revised retreat improve human-video transfer on unused seeds. These results suggest that the value of additional human data should be evaluated together with the transfer-error distribution it induces, since changes in timing, motion, or scene variation can shift the dominant transfer bottleneck.

\textbf{Author responsibility.} Reported values were checked against available experimental records; the authors are responsible for the work.
